% ============================================================================
% 城市级案例研究与实验分析 - 05_case_study.tex
% ============================================================================

\section{实验分析}
\label{sec:case_study}

xxxx

\subsection{真实城市地理环境与仿真数据源构建}

为验证本文算法在动态时空环境中的有效性，本研究需要构建具有真实物理与社会意义的城市时空风险场。由于真实的高精度动态人口/交通轨迹数据受限于隐私难以获取，本文基于城市真实 POI 分布底图，借鉴空间交互引力模型，构建了一套城市人口/交通潮汐密度模拟发生器。

具体而言，针对住宅区、写字楼与主干道等不同 POI 属性，分别赋予反相位的时域激活函数（如写字楼呈现典型的早晚双峰高斯分布 $\mathcal{N}(\mu, \sigma^2)$，住宅区则呈现日间低谷）。通过空间高斯核函数距离衰减叠加，生成了分辨率为 $\Delta t = 5$ 分钟、网格粒度为 $d_{res}$ 的全天候动态人口密度矩阵 $\mathbf{M}_{pop}(x,y,t)$ 作为本实验的底层数据支撑。详细的生成过程与数学推导参见附录 \ref{appendix:poi_model}。

\subsection{实验参数设置与基线算法说明（如与传统A*、单纯EDA等做对比）}



\subsection{极端时空场景的对比验证（论文高光时刻：证明你的模型对时间和天气极其敏感）}